%% kapa: exam=1 round=33 year=2022 period=1 form=A subject=civil_law subject_ko=민법 seq=32 number=32
%% source: Q-Net 감정평가사 기출문제 (cst00309) artlSeq=5209997 file=raw/1008/5209997/2022년도 제33회 감정평가사 1차시험 1교시(민법,경제학,부동산학).pdf
%% conversion: kordoc 4.24.51 -> tools/gen1.py (auto); review=proofread-sonnet; flags=none
\begin{question}{32}
전세권에 관한 설명으로 옳지 않은 것은? (다툼이 있으면 판례에 따름)
\begin{choices}
  \item 전세금의 지급은 전세권 성립의 요소이다.
  \item 기존 채권으로 전세금의 지급에 갈음할 수 있다.
  \item 농경지를 전세권의 목적으로 할 수 있다.
  \item 전세금이 경제사정의 변동으로 인하여 상당하지 아니하게 된 때에는 당사자는 장래에 대하여 그 증감을 청구할 수 있다.
  \item 전세권의 목적물의 전부 또는 일부가 전세권자에 책임있는 사유로 인하여 멸실된 경우, 전세권설정자는 전세권이 소멸된 후 전세금으로써 손해의 배상에 충당할 수 있다.
\end{choices}
\end{question}
\answer{3}
